\section{A line arrangement controls every ratio and dilation}
\label{sec:uniformity}

Fix a stable center $x$ and an atom $\xi$ with first default vertex $U$.
Write $r=r_h$ for its window length. Every local test $Q_i(z)$ consults
only the edge $e_i=(U,i)$ and its child subtree. Thus its outcome, for
\emph{every} realization of the remaining tables, is determined by
$J_{b_{e_i}^*}(z)$. We count the possible vectors
\begin{equation}\label{eq:key-vector}
 \left(J_{b_{e_i}^*}(x+ta_j(q)):
  1\le i\le m_0,\ j\in W_{e_i}\right).
\end{equation}

\subsection{Counting all faces, not only open regions}

\begin{lemma}[An elementary arrangement bound]\label{lem:line-arrangement}
For $L$ affine lines in the plane, the total number of vertices,
one-dimensional faces, and two-dimensional faces is at most $2L^2+1$.
After restriction to a rectangle, the number of sign vectors of $F$
affine functions is therefore at most $2(F+4)^2+1$.
The same bound holds if either pair of opposite sides is omitted, and
is also an upper bound on a vertical segment.
\end{lemma}

\begin{proof}
Discard duplicate lines. There are at most $L(L-1)/2$ intersection
points. Each line is divided by at most $L-1$ points into at most $L$
one-dimensional open pieces, giving at most $L^2$ edges. Adding the
lines one at a time gives at most $1+L(L+1)/2$ two-dimensional regions.
The sum is $2L^2+1$. Parallel lines, multiple intersections, and other
degeneracies only reduce these bounds.
For a rectangle, append its four supporting lines. Every face of the
resulting arrangement has constant signs, including zeros. Its
intersection with the required rectangle or open-sided rectangle is
empty or a single convex face. On a vertical segment the direct bound
$2F+3$ is smaller than the stated bound.
\end{proof}

The inclusion of zero signs is essential: a point exactly on a cell
boundary belongs to its right-hand cell. Sampling only the interiors
of two-dimensional regions would not justify the construction.

\begin{lemma}[Deterministic parameter representatives]\label{lem:representatives}
Let
\begin{equation}\label{eq:F-count}
 F_r=m_0r(1+4cQ^{-2r}).
\end{equation}
The vector \eqref{eq:key-vector} has at most
\begin{equation}\label{eq:D-count}
 D(r,B)=(2Tm_0rB+3)\bigl[2(F_r+4)^2+1\bigr]
\end{equation}
possible values. In particular,
\begin{equation}\label{eq:D-coarse}
 D(r,B)\le C(B+1)^4Q^{-4r},
\end{equation}
where the constant independent of $r_0,r,B$ is
\begin{equation}\label{eq:C-count}
 C=(2T+3)\bigl[2(5+4c)^2+1\bigr]m_0^3.
\end{equation}
One representative $(q,t)$ for each value can be chosen using only
$x$, the deterministic grids and windows, and $U$; the unexposed random
entries are irrelevant to this choice.
\end{lemma}

\begin{proof}
There are $m_0r$ relevant indices $j$, all at most $B$.
Partition $[\alpha,\beta]$ at the jumps of every relevant
$\nu_j(q)$. Retain each jump and endpoint as a singleton, and each
intervening open interval as a separate stratum. By
\cref{lem:envelopes}, at most $2Tm_0rB+3$ strata result.
On an open stratum every exponent $n_j=\nu_j(q)$ is constant.

For a pair $(i,j)$, the translated point lies in $(x,x+2Q^j]$.
Only boundaries $\gamma\in\Gamma_{b_{e_i}^*}\cap[x,x+2Q^j]$
can change its key. Their number is at most
\begin{equation}\label{eq:boundary-count}
 1+2N_{b_{e_i}^*}Q^j
 \le1+4cQ^{j-b_{e_i}^*}
 \le1+4cQ^{-2r},
\end{equation}
using \cref{lem:span}. A boundary with $\gamma\le x$ is never
attained by a translated point and may be discarded. For the others,
the boundary equation is
\[
 tq^{n_j}=\gamma-x.
\]
In coordinates
\begin{equation}\label{eq:log-coordinates}
 u=-\log q,\qquad v=\log t,
\end{equation}
it becomes the affine line
\begin{equation}\label{eq:boundary-line}
 v-n_ju=\log(\gamma-x).
\end{equation}
There are at most $F_r$ such lines, counting repetitions if necessary.
The logarithm is increasing, so its sign relative to a boundary is
exactly the sign of the original displacement relative to that boundary.
The sign vector, including its zeros, determines every key in
\eqref{eq:key-vector}.

On an open ratio stratum, \cref{lem:line-arrangement} supplies at most
$2(F_r+4)^2+1$ representatives. On a singleton stratum the exponents
are assigned by \eqref{eq:nu}, and a one-dimensional boundary count
gives the same upper bound. Multiplication by the number of strata
proves \eqref{eq:D-count}. All geometric data used here are fixed
under $\xi$.

Finally $r\le B$, $m_0\ge1$, and $Q^{-2r}\ge1$. Hence
\[
 2Tm_0rB+3\le(2T+3)m_0(B+1)^2,
 \qquad
 F_r+4\le m_0r(5+4c)Q^{-2r}.
\]
These inequalities give \eqref{eq:D-coarse} with \eqref{eq:C-count}.
\end{proof}

\begin{proposition}[Conditional uniform failure bound]\label{prop:uniform-failure}
For a stable $x$ and an atom $\xi$ whose first default has height $h$,
\begin{align}
 &\Prob\!\left(\exists(q,t)\in[\alpha,\beta]\times[1,2]
   \ \forall j\in\mathcal J:\ x+ta_j(q)\notin A
   \,\middle|\,\xi\right)\notag\\
 &\hspace{30mm}\le C(B+1)^4
  \bigl[Q^{-4}(1-p/2)^{m_0}\bigr]^{r_h}.
 \label{eq:uniform-failure}
\end{align}
\end{proposition}

\begin{proof}
If the selected points all miss $A$, \cref{lem:local-hit} implies
failure of all local tests at that parameter pair. A representative
from \cref{lem:representatives} has exactly the same key vector, so
all its local tests also fail for that realization of the tables.
Apply \cref{lem:failure} at each representative, and then the union
bound and \eqref{eq:D-coarse}. Because representatives were selected
before revealing the remaining entries, each application uses precisely
the fixed-parameter probability proved in \cref{lem:failure}.
\end{proof}

\subsection{Choose the parameters without a circular dependence}

Here is the complete order of choice.
First, after fixing $p,\alpha,\beta,Q,T,c$, choose $M$ so large that
\begin{equation}\label{eq:theta}
 \theta=Q^{-4}(1-p/2)^{M-1}<1.
\end{equation}
Second choose $d$ with $(1-2^{1-M})^d<p$.
Third compute $K$ from \eqref{eq:Kedges} and choose $g$ with
\begin{equation}\label{eq:gap-choice}
 4KcQ^{g+1}<p.
\end{equation}
The index $j_0$ was fixed using only $Q$.
Finally choose $r_0$ so large that
\begin{equation}\label{eq:base-length-choice}
 C(B(r_0)+1)^4\theta^{r_0}<p.
\end{equation}
Such a choice exists because \cref{lem:span} makes $B(r_0)$ affine
in $r_0$ while $0<\theta<1$. For every height $h$ we have
$r_h\ge r_0$, so \eqref{eq:uniform-failure} is less than $p$ with
this same choice. All comparisons in the choices can be made with
rational arithmetic when the initial parameters are rational.

Together with \eqref{eq:no-default}, this proves, for every stable $x$,
\begin{equation}\label{eq:pointwise-miss}
 \Prob\!\left(\exists(q,t)\in[\alpha,\beta]\times[1,2]
  \ \forall j\in\mathcal J:\ x+ta_j(q)\notin A\right)<2p.
\end{equation}
No claim is made that the required tree is of manageable size. The
parameter choice proves existence; its computational cost is discussed
in \cref{sec:effective}.

\subsection{Repair all centers using the original exponents}

We now prove \cref{prop:normalized}. For each table outcome, enlarge
the finite half-open cell union $A$ to an open periodic $A^+\supset A$
with $\rho(A^+)\le\rho(A)+p$. Define
\begin{equation}\label{eq:R-original}
 R=\left\{x\in\R:\exists(q,t)\in[\alpha,\beta]\times[1,2]
   \quad x+tq^n\notin A^+\ \text{for all }1\le n\le TB\right\}.
\end{equation}
The use of the \emph{original} exponents in this definition is deliberate.
The maps $q\mapsto a_j(q)$ jump. A projection argument applied directly
to those selected terms would not establish closedness.

By contrast, $x+tq^n$ is continuous for every original integer $n$.
On $(\R/\Z)\times[\alpha,\beta]\times[1,2]$ the joint miss set
in \eqref{eq:R-original} is closed. Its projection $R$ is therefore
closed and periodic. If a stable center belongs to $R$, then for the
witnessing $(q,t)$ all selected exponents also miss, because
$\nu_j(q)\le Tj\le TB$. Thus \eqref{eq:pointwise-miss} bounds
$\Prob(x\in R)$ by $2p$ on $\calS$. From
\eqref{eq:unstable-density} and \eqref{eq:gap-choice},
\begin{equation}\label{eq:expected-R}
 \E\rho(R)\le2p\rho(\calS)+\rho(\R\setminus\calS)\le3p.
\end{equation}
Also \cref{lem:mean-density} gives $\E\rho(A^+)\le2p$.
Choose an outcome with
\[
 \rho(A^+)+\rho(R)\le5p.
\]
Outer regularity on the circle supplies an open periodic $V\supset R$
with $\rho(V)\le\rho(R)+p$. Set $H=A^+\cup V$. Then
$\rho(H)\le6p$.

For $x\notin R$, every pair $(q,t)$ has a hit in $A^+$ among
$n=1,\ldots,TB$. For $x\in R$, openness of $V$ and convergence
$x+tq^n\to x$ give a hit in $V$ for all sufficiently large $n$.
This proves \eqref{eq:normalized-hit}, and hence
\cref{prop:normalized}. Notice that the repair indices need not lie
in the finite prefix used to define $R$.
